\section{Consultas obligatorias y cobertura}
Las siguientes tablas resumen la evidencia anual, las pruebas controladas y la
ruta local v2, diferenciadas por fila. Una implementación nueva se acredita con
su informe de ejecución, no solo con que el archivo exista.
``Demo'' significa entidades
administrativas sintéticas autorizadas; no inventario físico observado.
``Registrada'' no equivale a aceptación académica ni reproducibilidad actual.

\subsection{PostgreSQL}
\begin{longtable}{p{0.08\textwidth}p{0.39\textwidth}p{0.43\textwidth}}
\toprule ID & Consulta exigida & Evidencia y límite \\
\midrule\endhead
PG01 & Sensores por tipo. & Registrada sobre sensores demo. \\
PG02 & Sensores por ubicación. & Registrada sobre sensores demo y catálogo. \\
PG03 & Cantidad de activos e inactivos. & Registrada con estados administrativos demo. \\
PG04 & Sensores por organismo responsable. & Registrada con organismos asignados a sensores demo. \\
PG05 & Ubicaciones con más sensores. & Ranking registrado sobre sensores demo. \\
PG06 & Sensores instalados por año. & Registrada con fechas sintéticas. \\
PG07 & Sensores sin actividad registrada. & Funcional en prueba controlada; bloqueada en el lote observado por falta de relación sensor--canal. \\
PG08 & Ranking de zonas monitoreadas. & Registrada sobre sensores demo; revisar alcance geográfico. \\
PG09 & Consulta utilizando vistas. & Vista anual de actividad por canal. \\
PG10 & Procedimiento o función. & Función anual que cuenta canales de una estación. \\
\bottomrule
\end{longtable}

La regresión de integración fue corregida: el archivo
\archivo{sprint1/postgres/queries.sql} vuelve a contener PG01--PG10 para
\archivo{smartcity}; el núcleo local usa \archivo{staging_queries.sql}.
El parser anual reconoce las diez consultas y tiene pruebas de regresión.
Una nueva prueba PostgreSQL ejecuta el archivo canónico y contrasta las filas
de cada consulta con resultados esperados sobre fixtures revertidas al finalizar.

PG07 requiere una relación conocida sensor--canal antes de evaluar ausencia
de lecturas en enero. Excluye relaciones desconocidas y usa una ventana
semiabierta en hora de Madrid. La prueba incluye ambos límites, un sensor con
canal sin lecturas y un sensor sin canal; no acredita identidad física observada.

El núcleo demuestra SELECT, filtros, ordenamiento, DISTINCT, JOIN, GROUP BY,
HAVING, agregaciones, CTE, ventanas, subconsultas y fechas en distintos archivos.
Esa variedad técnica no reemplaza las preguntas obligatorias sobre sensores.

\subsection{MongoDB}
\begin{longtable}{p{0.08\textwidth}p{0.39\textwidth}p{0.43\textwidth}}
\toprule ID & Consulta exigida & Evidencia y límite \\
\midrule\endhead
MG01 & Lecturas de un sensor específico. & Búsqueda registrada por canal lógico, no dispositivo físico. \\
MG02 & Lecturas en un rango de fechas. & V2: rango BSON semiabierto, contrastado con conteo independiente del archivo. \\
MG03 & Máximos y mínimos. & V2: extremos de aire por variable/unidad, comparados con Python. \\
MG04 & Eventos críticos detectados. & V2: regla demo de intensidad superior a 2.000 veh/h; eventos derivados, no oficiales. \\
MG05 & Alertas emitidas por zona. & V2: emisión actual de replay en colección local; zonas del inventario disponible. \\
MG06 & Promedio de mediciones por zona. & V2: solo zonas conocidas; vigencia histórica no confirmada. \\
MG07 & Promedio por tipo de sensor. & V2: por fuente/variable/unidad; alternativa por tipo lógico, no físico. \\
MG08 & Incidentes por categoría. & V2: accidentes reales de 2026 deduplicados y anonimizados. \\
MG09 & Consulta con agregaciones. & V2: conteos por fuente y flag de calidad. \\
MG10 & Consulta con filtros múltiples. & V2: fecha, fuente, flag, valor y zona; resultado comparado con Python. \\
\bottomrule
\end{longtable}

El harness anual contiene algunas afirmaciones constantes en
\archivo{query_evidence}; su verificador controla estructura y procedencia,
no ejecuta cada pipeline. Por ello, no se declaran siete consultas anuales
completamente acreditadas solo por esas banderas. Agrupar una zona nula tampoco
demuestra un análisis por zona geográfica.

La ruta nueva no reutiliza esas banderas como prueba. El archivo
\archivo{sprint1/basico/consultas_mongodb_v2.js} contiene los diez pipelines y
filtros. \archivo{scripts/validar_v2.py} compara cantidades, extremos y promedios
con cálculos independientes sobre el JSONL verificado. Los informes nuevos
conservan parámetros, muestras y resultados agregados de cada consulta.

MG04 usa una regla académica de demostración, no un límite legal. MG05 guarda
\archivo{emitted_at} actual, \archivo{replay=true} y
\archivo{delivery=local_collection_only}; no fabrica alertas oficiales pasadas
ni envía notificaciones externas. El uso de distritos del inventario se declara
como \archivo{current_inventory_not_historically_verified}.

Las operaciones insertOne, insertMany, find, findOne, updateOne, updateMany y
deleteOne se prueban en la integración con colecciones de demostración aisladas.
Las actualizaciones y borrados no deben aplicarse a las observaciones originales.

\subsection{Ficha mínima para completar cada consulta}
Daniel y Kembly deben registrar, para cada ID:
\begin{enumerate}
\item Pregunta, contrato, esquema/colección y generación utilizada.
\item Archivo y SQL o pipeline exacto, con sus parámetros.
\item Resultado obtenido y comprobación independiente del valor esperado.
\item Índice/plan cuando corresponda, fecha, entorno y responsable de revisión.
\item Origen observado, derivado o sintético y limitaciones de interpretación.
\end{enumerate}
La matriz editable del repositorio está en
\archivo{docs/equipo/MATRIZ_CONSULTAS_SPRINT1.md}.
